% Pursuing Stacks (cotes 134-2 to 134-8): archive cote 134-2, site page 15 (PDF page 16).
% letter to Quillen. Draft compared: ps1.tex:52-108.
% AI-transcribed by Claude Opus 5.5 (claude-opus-5-5, adaptive thinking, effort high), Message Batches API, batch ps1 (msgbatch_01Rv5hCaNTtUt4JHx1yqoHPL), 2026-11-12; not yet checked by eye.
\scanpage{134-2/15}
% NOTE: the upper third of the scan carries no text on this page; only mirror-image show-through from another sheet is visible there ("Chapter I : TAKE-OFF", "A letter to ....", "(followed by an appendix: Three letters to Larry Breen)", and a faint "2'"); it is not transcribed
% NOTE: the next two lines are written by hand over a strip of paper; faint typed text shows beneath them and cannot be read
\hand{any program of homological, homotopical or topological algebra}
\margin{\unsure{=}}\hand{2. The motivation for this letter was the latter topic however. Ronnie}
Brown and his friends are competent algebraists and apparently strongly
motivated for investing energy in foundational work,  on the other hand
they visibly are lacking the necessary scope of vision which geometry
alone provides\hand{$^{(3)}$}.They seem to me kind of isolated,  partly due I guess to
the disrepute I mentioned before  -  I suggested to try and have contact
with people such as yourself, Larry Breen, Illusie and others,  who have
the geometric insight  and \unsure{who\xout{m}|whom} moreover, \xout{\illegible{1}} may not think themsel-%
ves to\inserted{o} \unsure{ggod|good} for indulging in occasional reflection o\over{f}{n} foundational matters\hand{\unsure{,}}
and in the process help others do the work which should be done.

At first sight it had seemed to me that the Bangor group\hand{\unsure{$^{(5)}$|$^{(4)}$}}\hand{/}had indeed
come to work out (quite independently) one basic intuition of the program
I had envisioned in those letters to Larry Breen  -  namely that the
study of n-truncated homotopy types (of semisimplicial sets,  or of topo-%
logical spaces) was essentially equivalent to the study of so-called
n-groupoids (where n is any natural integer)\unsure{.} This is ex\hand{p}ected to be
achieved by associating to any space (say) X it's\hand{|}``fundamental n-groupoid''
\hand{$\Pi$}$_n$(X),  generalizing the familiar Poincar\'e fundamental groupoid for n=1.
The \unsure{\over{a}{o}bvious|obvious} idea is that 0-objects of \hand{$\Pi$}$_n$(X) should be the points of X,
1-objects should be ``homotopies'' or paths between points,  2-objects
should be homotopies between 1-objects, etc.  This \hand{$\Pi$}$_n$(X)  should embody
the n-truncated homotopy type of X, in much the same way as for n=1 the
usual fundamental groupoid embodies the 1-truncated homotopy type.  For
two spaces X,Y,  the set of homotopy-classes of maps X\hand{$\to$}Y (more correct-%
ly, for general X,Y,  the maps of X into Y in the homotopy category)
should correspond to n-equivalence classes of n-functors from \hand{$\Pi$}$_n$(X) to
\hand{$\Pi$}$_n$(Y)  -  etc.  There are very stro\hand{$\frown$}ng  suggestions for a nice formalism,
including a notion of geometric realization of an n-groupoid,  which
should imply that any n-groupoid is n-equivalent to a \hand{$\Pi$}$_n$(X).  Moreover,
when the notion of n-groupoid (or more generally of an n-categroy) is
relativized over an arbitrary topos to the notion of an n-gerb (or more
generally, an n-stack), these become the \xout{\illegible{1}} natural ``coefficients'' for
a formalism of non commutative cohomological algebra,  in the spirit of
Giraud's thesis.

But all this kind of thing for the time being is pure heuristics
-  I never so far sat down to try to make explicit at least a definition
of n-categories and n-groupoids,  of n-functors between these etc.  When
I got the Bangor reprints I at once had the feeling that this kind of \hand{\unsure{work|with}}

\mspage{15}
